\section{System}
\label{sec:system}

The One runs teams of language-model agents on a person's computer or on a Linux server. The
desktop application and the headless server, \code{the-one-server}, share one backend; every
experiment in this paper ran the server \citep{zeroth2026server}.

\para{Agents with a life loop} An agent is a durable record in SQLite, not a process. Its state is
committed before anything acts on it, and a woken agent continues from that record, so a crash or
restart loses no decided work. A turn ends only through a lifecycle tool (deliver, hand back,
wait); a model reply without one continues the loop, and errors are retried rather than ending the
run.

\para{Teams as graphs} A team is a YAML definition: agents with roles, prompts, tools and model
bindings, wired as a graph of who may hand work to whom. A lead can open a subgraph (a clone of
itself or a helper team), hand work to another node, or ask the person. The configuration we kept,
\emph{team 7.0}, is a single lead that works directly and delegates only when asked; the earlier
shipped teams (6.0 to 6.6) added checker nodes and rules, described in \cref{sec:structure}.

\begin{figure}[t]
\centering
\begin{tikzpicture}[
  font=\small,
  node distance=7mm and 14mm,
  box/.style={draw=ink!70, rounded corners=2pt, align=center, inner sep=3pt,
              minimum height=11.5mm, text width=24mm, fill=white},
  ext/.style={box, fill=boxfill},
  arr/.style={-{Stealth[length=2.1mm]}, semithick, draw=ink},
  darr/.style={{Stealth[length=2.1mm]}-{Stealth[length=2.1mm]}, semithick, draw=ink},
  lab/.style={font=\scriptsize, text=ink!85, align=center, inner sep=1.5pt}
]
  \node[ext] (harbor) {\textbf{Harbor}\\{\scriptsize one container}\\{\scriptsize per attempt}};
  \node[box, right=of harbor] (install) {installs the\\published server};
  \node[box, right=of install] (run) {\code{the-one-server}\\\code{run}\\{\scriptsize instruction on stdin}};
  \node[box, below=of run] (lead) {\textbf{lead agent}\\{\scriptsize team 7.0}};
  \node[box] (folder) at (install |- lead) {task folder\\{\scriptsize shell, files}};
  \node[ext, right=of lead] (model) {\textbf{model API}\\{\scriptsize GPT-6.2 Sol, max}};
  \node[box, below=of lead] (out) {\code{run.json},\\ATIF trajectories};
  \node[box] (verify) at (install |- out) {verifier\\{\scriptsize hidden tests}};
  \node[ext] (reward) at (harbor |- out) {\textbf{reward}\\{\scriptsize 0 or 1}};

  \draw[arr] (harbor) -- (install);
  \draw[arr] (install) -- (run);
  \draw[arr] (run) -- (lead);
  \draw[darr] (lead) -- node[lab, above] {model\\calls} (model);
  \draw[darr] (lead) -- node[lab, above] {tool\\calls} (folder);
  \draw[arr] (lead) -- node[lab, right, align=left, xshift=1.5pt] {delivers or\\times out} (out);
  \draw[arr] (out) -- node[lab, above] {then} (verify);
  \draw[arr] (verify) -- (reward);

  \begin{scope}[on background layer]
    \node[draw=ruleGray, dashed, rounded corners=4pt, inner sep=3.2mm,
          fit=(install) (run) (lead) (folder) (out) (verify)] (ctr) {};
  \end{scope}
  \node[lab, anchor=south west] at (ctr.north west) {task container};
\end{tikzpicture}
\caption{One Terminal-Bench attempt through the Harbor adapter, as run for the final evaluation.
In team 7.0 the lead does the work itself and opens helpers only when asked; every other box is the
same for any agent run through Harbor. The agent's own record (\code{run.json}) holds the outcome,
token usage and list-price cost; the reward comes only from the verifier.}
\label{fig:flow}
\end{figure}

\para{Context epochs} A model request is a stable base (instructions, skills, project notes), an
append-only log, and a short tail (\cref{fig:epochs}). New context is appended where it appears
and then stays put, so the provider's prompt cache keeps hitting. The base is re-rendered only at
named boundaries: a change of model or tools, a new summary, or more than 40k tokens of drift. The
earlier (``classic'') compiler re-rendered the instructions on every cycle.

\begin{figure}[t]
\centering
\begin{tikzpicture}[font=\small, x=1mm, y=1mm]
  % classic layout
  \node[anchor=east, align=right] at (-2,14) {classic\\[-2pt]{\scriptsize each cycle}};
  \fill[oursLight!70] (0,10) rectangle (38,18);
  \node at (19,14) {instructions, re-rendered};
  \fill[gray!16] (38,10) rectangle (96,18);
  \node at (67,14) {history};
  \fill[gray!34] (96,10) rectangle (112,18);
  \node at (104,14) {tail};
  \draw[line width=1.8pt, ours] (0,7.5) -- (5,7.5);
  \node[anchor=west, font=\scriptsize, text=ours!85!black] at (6.5,7.4)
    {cache often breaks inside the system part (30--40\% of requests)};
  % epoch layout
  \node[anchor=east, align=right] at (-2,-6) {epochs};
  \fill[gray!30] (0,-10) rectangle (38,-2);
  \node at (19,-6) {stable base};
  \fill[gray!16] (38,-10) rectangle (96,-2);
  \node at (67,-6) {append-only log};
  \fill[gray!34] (96,-10) rectangle (112,-2);
  \node at (104,-6) {tail};
  \draw[line width=1.8pt, ink!75] (0,-12.5) -- (96,-12.5);
  \node[anchor=north, font=\scriptsize] at (48,-13.5)
    {cached prefix: 89--94\% of prompt tokens read from cache (paired run, 20 tasks)};
\end{tikzpicture}
\caption{Request layout under the classic compiler and under context epochs. Epochs keep the base
fixed between named boundaries (model or tool change, new summary, more than 40k tokens of drift)
and only append, so everything before the tail stays a cache hit.}
\label{fig:epochs}
\end{figure}

\para{Summaries} When a history grows long, a summarizer (``meditation'') replaces it with a
compact summary. The working assignment, its amendments and the newest report are carried
verbatim, because summaries otherwise kept only about a third of rule sentences
(\cref{sec:context}).

\para{Headless runs and the Harbor adapter} \code{the-one-server run} takes a task on standard
input, runs the team in the task's folder with nobody to answer its questions, and writes
\code{run.json} (outcome, usage, list-price cost, changed files) and an ATIF trajectory for every
agent. A Harbor installed agent installs the published server in each task container and calls
\code{run} (\cref{fig:flow}). The adapter is public \citep{zhang2026adapter}; the command that
reproduces the final run is in \cref{app:repro}.
